case where the base is a trait. It follows that the properties of the specialization homomorphism reviewed in the preceding sections, in the case where the base is a trait, extend automatically to the case of a more general base. Of course, the methods of definition considered apply equally well to @@M0@@, @@M1@@ (when @@M2@@ is smooth), and @@M3@@ (when, in addition, @@M4@@ is separated).

\sgaStableTarget{sga6.e10.para.7-17-7}\textbf{7.17.7.} Finally, note that the definition of a homomorphism (7.17.3.1) or (7.17.3.2), or of a homomorphism (7.17.4.2), depends essentially on certain arbitrary choices (a suitable trait over @@M5@@ in 7.17.3, a sequence of @@M6@@ in 7.17.4, and, in addition, a resolution of singularities in 7.17.5), and likewise for the variants for @@M7@@, etc. Obvious examples can be constructed with @@M8@@, say, an ordinary double point of an algebraic curve and @@M9@@ a projective smooth scheme over @@M10@@ (for example, an abelian scheme of relative dimension 2), in which the specialization homomorphism on @@M11@@ and even on @@M12@@ (the Néron--Sévéri group) actually depends on the choice of the trait @@M13@@ chosen over @@M14@@ or, what amounts to the same thing here, on the choice of a point over @@M15@@ in the normalization of @@M16@@.

\sgaStableTarget{sga6.e10.sec.7.18}\subsection*{7.18. Appendix to the Appendix: extension of the base field for certain functors with values in commutative groups}
We spell out here a technical point raised in 7.17.4.

\sgaStableTarget{sga6.e10.para.7-18-1}\textbf{7.18.1.} Let @@M17@@ be a field, @@M18@@ the category of integral @@M19@@-algebras and flat homomorphisms of such algebras, and @@M20@@ a functor on @@M21@@ with values in the category @@M22@@ of sets. Suppose that @@M23@@ commutes with filtered inductive limits. Let @@M24@@ be an extension of @@M25@@; we seek
